\section{The complete equality language}
\label{sec:equality}

The potential proves more than nonnegativity. Equality in
\eqref{eq:exact-endpoint-sum} forces every transformed weight along the
path to vanish, and this reduces an inequality over $11420$ states to a
language on six of them.

\begin{theorem}[Equality and the reflected word]
\label{thm:equality-language}
Under the hypotheses of \cref{thm:endpoint-89}, equality $\Delta=0$
holds exactly when
$$
 q=111\chi(s)2,\qquad s\in(12111\mid21111)^*.
$$
The target is $111\chi(\iota(s))2$, where $\iota$ exchanges the two
five-letter tokens, preserving their order in the token sequence.
\end{theorem}

\subsection{Exhausting the zero-slack paths}

Retain every edge whose transformed weight is zero. In this graph, keep
the states that are both reachable from a zero-slack seed and able to reach
a zero-slack primitive terminal. The complete calculation gives exactly
the states in \cref{tab:equality-states}. The only surviving seed is
branch $R_0$, optional digit $1$, with empty paired prefix and state $s_0$.
The only surviving terminal is $s_1$.

\begin{table}[htbp]
\centering
\caption{The full equality graph. The last column is the terminal charge
$f(S)$ when the state is a primitive terminal; a dash means the state is
not such a terminal. Only $s_0$ is an initial state and only $s_1$ is accepting.}
\label{tab:equality-states}
\begin{tabular}{clrc}
\toprule
State & Matrix $S$ & $\lambda(S)$ & $f(S)$\\
\midrule
$s_0$ & $(119,39;41,80)$ & $-1$ & --\\
$s_1$ & $(115,37;47,84)$ & $-1$ & $1$\\
$s_2$ & $(89,21;0,89)$ & $0$ & --\\
$s_3$ & $(89,0;68,89)$ & $-1$ & --\\
$s_4$ & $(131,21;47,68)$ & $-1$ & --\\
$s_5$ & $(121,37;41,78)$ & $-1$ & --\\
\bottomrule
\end{tabular}
\end{table}

Every edge is shown in \cref{fig:equality-graph}. From $s_0$, the first
letter must be $1$, reaching $s_1$. The two first-return words at $s_1$
are exactly $12111$ and $21111$. Thus the accepted paired words are
$1(12111\mid21111)^*$. The optional initial digit $1$ converts this
to the asserted source language, which proves necessity.

That the returns at $s_1$ are two in number, and not one or three, is what
makes the equality language a free monoid on two generators, and hence what
makes the count relation $Q=5P$ of \cref{cor:genuine-strict} rigid. The
same numerical signature characterizes Sturmian words, whose every factor
has exactly two return words \cite{Vuillon2001,JustinVuillon2000}. The two notions are
distinct: returns to a factor of an infinite word on one side, first
returns to a state of a weighted graph on the other. We assert no
implication between them. Both tokens above nevertheless carry one heavy
letter in five, which is the frequency $\theta=1/5$ that these same six
states contribute to the critical set of \cref{sec:markoff-application}.

\begin{figure}[tb]
\centering
\begin{tikzpicture}[>=Stealth,
 state/.style={circle,draw,minimum size=8mm,inner sep=1pt},
 every edge/.style={draw,->},every node/.style={font=\small}]
\node[state] (s0) at (0,0) {$s_0$};
\node[state,double,double distance=1.2pt] (s1) at (2.1,0) {$s_1$};
\node[state] (s2) at (4.1,1.15) {$s_2$};
\node[state] (s3) at (4.1,-1.15) {$s_3$};
\node[state] (s4) at (6.2,0) {$s_4$};
\node[state] (s5) at (3.1,-2.7) {$s_5$};
\draw[->] (-1,0)--(s0);
\path (s0) edge node[above] {$1:0$} (s1)
 (s1) edge node[above left] {$1:1$} (s2)
 (s1) edge node[below left] {$2:0$} (s3)
 (s2) edge node[above right] {$2:-1$} (s4)
 (s3) edge node[below right] {$1:0$} (s4)
 (s5) edge[bend left=18] node[below left] {$1:0$} (s0);
\draw[->] (s4.east) .. controls (8,0) and (8,-2.7) ..
 node[right,pos=.52] {$1:0$} (s5.east);
\end{tikzpicture}

\caption{All zero-slack accepted paths at norm $89$. An arrow label
$j:w$ gives the paired input letter and its original edge weight;
the transformed weight is zero on every displayed arrow.
The seed is $s_0$ and the double circle marks the only accepting state
$s_1$. Its two first-return words are $12111$ and $21111$.}
\label{fig:equality-graph}
\end{figure}

\subsection{An exact token exchange}

The literal reflected word follows from small matrix identities. Put
\begin{equation}
 Q_0=M(1)^3,\qquad
 K=\begin{pmatrix}47&84\\68&-47\end{pmatrix},\qquad
 H=\begin{pmatrix}89&68\\0&-89\end{pmatrix}.
 \label{eq:token-intertwiners}
\end{equation}
Direct multiplication gives
\begin{equation}
 R_0Q_0=Q_0K,\qquad KM(2)=M(2)H.
 \label{eq:boundary-intertwiners}
\end{equation}
Writing $G_t=M(\chi(t))$, the two return identities are
\begin{equation}
 KG_{12111}=G_{21111}K,\qquad
 KG_{21111}=G_{12111}K.
 \label{eq:token-exchange}
\end{equation}
Induction over the token sequence therefore yields
\begin{equation}
 R_0M(111\chi(s)2)=M(111\chi(\iota(s))2)H.
 \label{eq:full-token-identity}
\end{equation}
The first column of $H$ is $89e_1$. Hence the reflected first column
is $89$ times the primitive positive continuant column of the displayed
target. The source and target have equal digit sums. This proves
sufficiency and the target formula in \cref{thm:equality-language},
including exact primitive acceptance, for every token sequence.

\begin{corollary}[Strict cost increase for genuine sources]
\label{cor:genuine-strict}
Let $v$ be the canonical source of a reduced Markoff occurrence, with
$\ord_{89}(\|v\|^2)=1$. Every distinct complete norm-$89$ target has cost
at least $\ell(v)+1$. The only genuine equality source is
$q=1112$, with $v=(8,5)^{\mathsf T}$, which the flip fixes.
\end{corollary}

\begin{proof}
If an equality word contains $j$ five-letter tokens, its inferred counts
are
$$
 P=1+2j,\qquad Q=5+10j=5P.
$$
A genuine reduced occurrence has $\gcd(P,Q)=1$ by definition, and
\cref{prop:genuine-source-dictionary} identifies the counts of its word
with that $P$ and $Q$. Therefore $j=0$. The base column is $(8,5)^{\mathsf T}$
and \eqref{eq:full-token-identity} fixes it after division by $89$.
All other genuine sources have positive integral excess.
\end{proof}

The proof uses coprime counts only after exhausting the equality
language. It does not use the converse assertion that a balanced paired
word with coprime counts is automatically a genuine occurrence.
